% ─────────────────────────────────────────────────────────────────────────────
\section{Engineering Impact, Ethics, and Contemporary Issues}
\label{sec:impact}
% ─────────────────────────────────────────────────────────────────────────────

The project is not only a technical exercise in computer vision. It also has broader global,
economic, environmental, societal, and ethical implications. Because the application domain
touches surveillance and security, these considerations are part of the design itself rather than
an afterthought.

\subsection{Global and Strategic Impact}

At a global level, the project addresses a widely relevant engineering problem: how to obtain
useful situational awareness from difficult imagery without depending on remote cloud
infrastructure. This matters not only in defense settings, but also in border monitoring,
critical-infrastructure protection, disaster assessment, and emergency response in degraded
visual environments.

The architectural emphasis on local processing and structured interpretation also contributes to
the broader trend toward edge intelligence. Instead of sending raw mission imagery to a
remote platform, the system tries to extract meaning close to the data source.

\subsection{Economic Impact}

The project was intentionally designed around open-source models and frameworks. This has
clear economic consequences:

\begin{itemize}[nosep]
  \item lower licensing cost compared with proprietary end-to-end platforms,
  \item lower recurring bandwidth and cloud-inference cost due to local processing,
  \item easier technology transfer to future academic or industrial teams because the toolchain
    is based on widely available frameworks.
\end{itemize}

At the same time, the project also demonstrates that ``open-source'' does not mean ``free of
engineering cost.'' Thermal adaptation, fine-tuning, validation, and deployment all still
require time, hardware, and expert effort.

\subsection{Environmental Impact}

The environmental profile of the system is mixed. Local edge inference can reduce the energy
and infrastructure cost associated with constant remote data transfer and cloud processing.
However, multimodal reasoning models are computationally expensive, and their training or
adaptation can still consume substantial compute resources.

This is why optimization techniques such as parameter-efficient fine-tuning and
quantization are important. They are not only performance strategies, but also part of a more
responsible use of computing resources.

\subsection{Societal and Ethical Impact}

The potential societal benefit of the project is improved operator awareness. A system that
helps distinguish between ``an object exists'' and ``this object is likely relevant or risky'' can
reduce fatigue and help prioritize attention in time-sensitive environments.

At the same time, surveillance technology creates ethical risk. A system that interprets people,
vehicles, and possible threats must not be treated as an unquestionable authority. The project
therefore follows a human-in-the-loop philosophy consistent with professional engineering
ethics~\cite{acm,ieeeethics}:

\begin{itemize}[nosep]
  \item the system informs human decision-making rather than replacing it,
  \item outputs are treated as analytical suggestions, not final truth,
  \item privacy, security, and data sovereignty are design requirements, not optional features.
\end{itemize}

\subsection{Contemporary Issues Related to the Project Topic}

Several contemporary issues are directly relevant to the project.

\begin{itemize}[nosep]
  \item \textbf{AI in security and defense:} there is growing pressure to increase automation,
    but this increases the importance of accountability and operator oversight.
  \item \textbf{Bias in multimodal models:} general VLMs are shaped by RGB-heavy training
    corpora, which can lead to systematic misunderstanding of thermal scenes.
  \item \textbf{Data sovereignty:} many current AI systems assume cloud-based deployment,
    which is unacceptable in some operational environments.
  \item \textbf{Open-source model governance:} open models improve accessibility and
    reproducibility, but they also require careful evaluation before domain-sensitive use.
  \item \textbf{Explainability vs. fluency:} a model can produce convincing language that is not
    always well grounded in thermal evidence, so output structure and operator review remain
    essential.
\end{itemize}

\subsection{Professional Responsibility}

From an engineering ethics perspective, the project reinforces three responsibilities:

\begin{itemize}[nosep]
  \item build systems that reduce harm rather than simply increasing technical capability,
  \item preserve transparency about what the prototype can and cannot do,
  \item treat validation, documentation, and human oversight as core deliverables.
\end{itemize}

For this reason, the final report does not present the system as an autonomous authority. It is
more accurate and more responsible to describe it as a secure, locally deployable semantic
assistant for thermal imagery.
